\documentclass[10pt]{article}
\usepackage{icmlko}

\icmlmeta{IP 파운데이션 월드모델}{205}{열 하나가 도메인 하나를 죽이고 있었다}{트리에서도 풀이 압축하나 물었다. 압축은 오히려 더 심한데 빼야 할 도메인이 바뀐다 --- 그리고 그 까닭이 열 하나로 좁혀졌다}

\begin{document}

\twocolumn[
\icmlheader

\begin{abstract}
노트 204가 능형에서 ``풀링 $=$ 0.127 $+$ 0.688 $\times$ 혼자''를 재고 ``트리는
도메인으로 쪼갤 수 있으니 덜 압축할 것''을 걸었다. \textbf{반대로 틀렸다} ---
트리는 절편 \textbf{0.207} · 기울기 \textbf{0.567} 로 능형(0.148 · 0.648)보다
\emph{더} 압축한다. 걸어 둔 반대 가능성(``트리 혼자 모형이 작은 도메인에서
과적합한다'')이 맞았다 --- 도서(학습 80행)에서 트리 혼자가 $+$0.116 으로
능형($+$0.290)보다 훨씬 나쁘다. 그런데 \textbf{빼야 할 도메인이 바뀐다} ---
능형은 모바일(격차 $+$0.075), 트리는 \textbf{웹툰}($+$0.077)이다. 그 뒤집힘이
이 노트의 본론이다. \textbf{웹툰 혼자 학습이 능형 $+$0.0371 인데 트리
$+$0.4495 다}(차 $+$0.41). 노트 182가 ``웹툰은 자기 과거를 못 쓴다''를
자기 전이 $+$0.037 로 세웠는데, \textbf{그 수는 맞고 그 문장은 넓다} ---
못 쓰는 것이 아니라 \emph{선형으로} 못 쓴다. 트리의 혼자 우위는 여섯
도메인에서 나타나고 학습 표본과 $\rho{=}+0.857$ 로 붙는다. 그다음 왜인지
좁혔다. \textbf{열 하나다.} \texttt{goods\_scale}(주당 회차 수) 하나를 빼면
웹툰 혼자 능형이 $+$0.0902 에서 \textbf{$+$0.3250} 으로 뛴다. 그 열은
능형의 \emph{최대} 계수($+$0.267)이고 노트 182가 잰 대로 학습 $+$0.378 에서
유보 $-$0.283 으로 부호가 뒤집히는 열이다. \textbf{능형은 그 열을 못
버리고 트리는 덜 쓴다.} 그리고 짝을 이루는 반례도 있다 ---
\texttt{target\_breadth} 도 시드는데 빼면 오히려 $-$0.0659 로 나빠진다.
\textbf{시든다고 다 버릴 것이 아니다.} 규제로는 절반만 된다 --- $\alpha$ 를
1 에서 500 으로 올리면 $+$0.037 에서 $+$0.264 인데, 열 하나 빼는 것($+$0.325)이
더 낫다.
\end{abstract}
]

\begin{figure*}[t]
\centering
\includegraphics[width=\textwidth]{figs/onecol.pdf}
\caption{왼쪽 --- 트리가 더 압축한다. 가운데 --- 웹툰만 $+$0.41.
오른쪽 --- 굿즈 규모 하나를 빼면.}
\end{figure*}

\section{트리가 더 압축한다}

\begin{table}[t]
\centering\tsize
\caption{풀링 $= a + b \times$ 혼자.}
\begin{tabular}{lrrrr}
\toprule
& 절편 $a$ & 기울기 $b$ & $R^2$ & 손익분기 \\
\midrule
F6 능형 & .148 & \textbf{.648} & .783 & .421 \\
\textbf{F18 트리} & \textbf{.207} & \textbf{.567} & .529 & .478 \\
\bottomrule
\end{tabular}
\end{table}

\claimx{\textbf{예측이 반대로 틀렸다.} ``트리는 도메인으로 쪼갤 수 있으니
덜 압축할 것''으로 기울기 0.85$\sim$1.0 을 적었는데 0.567 이다.}{걸어 둔
반대 가능성이 맞았다 --- 트리 혼자 모형이 작은 도메인에서 무너진다. 도서
(학습 80행)가 트리 $+$0.116 대 능형 $+$0.290 이다. \textbf{약한 쪽이 더
약해지니 업어 주는 폭이 커진다.}}

\begin{table}[t]
\centering\tsize
\caption{빼야 할 도메인이 바뀐다(격차 $=$ 혼자 $-$ 풀링).}
\begin{tabular}{lrr}
\toprule
& F6 능형 & F18 트리 \\
\midrule
\textbf{모바일} & \textbf{$+$.075} & $+$.007 \\
\textbf{웹툰} & $-$.202 & \textbf{$+$.077} \\
게임 & $-$.060 & $-$.061 \\
도서 & $-$.008 & \textbf{$-$.305} \\
\bottomrule
\end{tabular}
\end{table}

\claimx{\textbf{웹툰이 뒤집힌다} --- 능형에서는 풀에 꼭 넣어야 하고
($-$0.202) 트리에서는 빼는 것이 낫다($+$0.077).}{\textbf{노트 197$\sim$204 의
답이 형태마다 다르다는 뜻이다.} 그 갈래는 능형으로만 돌았고, 채택한 이분
규칙도 능형 이야기다.}

\section{웹툰은 선형으로 못 쓴다}

\begin{table}[t]
\centering\tsize
\caption{도메인마다 혼자 학습. 씨앗 셋 평균(트리).}
\begin{tabular}{lrrrr}
\toprule
& 능형 & 트리 & 차 & 학습 n \\
\midrule
\textbf{웹툰} & \textbf{.037} & \textbf{.450} & \textbf{$+$.413} & 2{,}106 \\
세계애니 & .469 & .542 & $+$.073 & 2{,}648 \\
모바일 & .507 & .571 & $+$.063 & 1{,}559 \\
펀딩 & .171 & .224 & $+$.053 & 320 \\
게임 & .398 & .441 & $+$.043 & 259 \\
애니 & .450 & .481 & $+$.031 & 1{,}467 \\
\textbf{도서} & .290 & .116 & \textbf{$-$.174} & \textbf{80} \\
\bottomrule
\end{tabular}
\end{table}

\claimx{\textbf{노트 182의 문장을 좁힌다.} ``웹툰은 자기 과거를 못 쓴다''는
자기 전이 $+$0.037 로 세운 것인데, 그것은 \emph{능형} 자기 전이다. 트리로는
$+$0.450 이다.}{수는 정확히 재현된다(여기 능형 $+$0.0371). \textbf{틀린 것은
측정이 아니라 일반화다} --- 이 흐름에서 여덟 번째다.}

\claimx{\textbf{트리 우위가 학습 표본과 붙는다}($\rho{=}+0.857$). 도서 80행
에서만 트리가 진다.}{쪼갤 재료가 있어야 쪼갠다 --- \textbf{작은 도메인에는
선형이 낫다.}}

\section{열 하나}

\begin{table}[t]
\centering\tsize
\caption{웹툰 혼자 능형. 학습 2,106 · 유보 711.}
\begin{tabular}{lr}
\toprule
전체 축 · $\alpha{=}1$ & .0371 \\
전체 축 · $\alpha{=}20$ & .0902 \\
전체 축 · $\alpha{=}500$ & .2643 \\
\midrule
\texttt{target\_breadth} 뺌 & \textbf{$-$.0659} \\
시든 둘 뺌 & .2979 \\
\textbf{\texttt{goods\_scale} 뺌} & \textbf{.3250} \\
\midrule
(트리) & .4495 \\
\bottomrule
\end{tabular}
\end{table}

\claimx{\textbf{\texttt{goods\_scale} 하나를 빼면 $+$0.0902 가 $+$0.3250 이
된다.} 그 열은 능형의 최대 계수($+$0.267)이고, 노트 182가 잰 대로 학습
$+$0.378 에서 유보 $-$0.283 으로 부호가 뒤집힌다.}{\textbf{능형은 그 열을
못 버린다} --- 학습에서 제일 강한 신호라 계수를 크게 준다. 트리는 다른
열로 쪼갤 수 있어 덜 쓴다.}

\claimx{\textbf{그런데 \texttt{target\_breadth} 는 빼면 나빠진다}($-$0.0659).
그 열도 학습 $+$0.319 에서 유보 $+$0.235 로 시드는데(노트 182) 부호는 안
뒤집힌다.}{\textbf{시든다고 다 버릴 것이 아니다} --- 노트 182가 시들기를
가드로 만들려다 실패한 자리(문턱 한 칸에 부호가 뒤집힘)의 까닭이 여기
있다. \textbf{버려야 하는 것은 시드는 열이 아니라 \emph{부호가 뒤집히는}
열이다.}}

\claimx{규제로는 절반만 된다 --- $\alpha$ 를 1 에서 500 으로 올리면
$+$0.037 이 $+$0.264 인데 열 하나 빼는 것($+$0.325)이 더 낫다.}{규제는 모든
계수를 같이 줄이므로 쓸모 있는 열도 같이 줄인다. \textbf{한 열만 틀렸을
때는 그 열만 빼는 것이 맞다.}}

\begin{table}[t]
\centering\tsize
\caption{남은 것.}
\begin{tabular}{ll}
\toprule
\textbf{부호 뒤집힘} & 학습 안에서 미리 잡을 수 있나 \\
트리 잔여 & 열을 빼도 .325 대 .450 --- 나머지 .125 \\
갈래 범위 & 노트 197$\sim$204 를 트리로 다시 \\
\bottomrule
\end{tabular}
\end{table}

첫 줄이 곧바로 잴 수 있고 값이 크다. 노트 182가 ``시들기''로 가드를 만들려다
문턱이 불안정해 실패했는데, 이 노트가 그 까닭을 준다 --- \textbf{시드는
것과 뒤집히는 것을 안 갈랐다.} \texttt{goods\_scale} 은 뒤집히고
\texttt{target\_breadth} 는 시들기만 한다. \textbf{학습 구간 안에서 연도별
상관이 0 을 가로지르는지 보면 둘을 가를 수 있고}, 그러면 유보를 안 보고
``버려야 할 열''을 짚을 수 있다. 노트 182의 최근 3해 기울기는 크기만 봤고
0 을 가로지르는지는 안 봤다.

\end{document}
